\section*{NS100: geometry reassessment and a complete theta-slice control}
\paragraph{Status.} This is an independent route assessment with a bounded
control, not a new RH criterion, a cofinal arithmetic bound, or a manuscript
version. Lagarias's half-plane logarithmic-derivative criterion and the
Jensen--P\'olya/Csordas associated-kernel program are prior art. No global
novelty is claimed for the control or the proposed tests.

\paragraph{A geometric observable that includes the function.}
Let $F(z)=\xi(1/2+z)/2$ and, where $F\ne0$, let $L=F'/F$.
The Cayley map $S=(L-1)/(L+1)$ satisfies
\[
 1-|S|^2=\frac{4\Re L}{|L+1|^2}.
\]
Thus a disk picture can encode a genuine analytic positivity question;
the picture does not prove the sign. The classical condition $\Re L(z)>0$
throughout $\Re z>0$ is equivalent to RH. It includes absence of poles
there. Merely drawing the image circle or imposing the functional equation
supplies none of that condition. NS94's count/similarity control remains valid.

\paragraph{Keep the full theta integral.}
In the normalization used here, for $t\ge0$ put
\[
 \phi(t)=\sum_{n\ge1}(2\pi^2n^4e^{9t/2}-3\pi n^2e^{5t/2})
                     e^{-\pi n^2e^{2t}},\qquad\phi(-t)=\phi(t).
\]
The classical theta identity gives $F(z)=\int_{\mathbb R}\phi(t)e^{zt}\,dt$.
The even extension is the complete theta kernel, not termwise evenness
of the individual exponentials. All derivatives and integrals below
converge on compact $z$ sets by its double-exponential real tail.
For $x>0$, $y\in\mathbb R$, set
\[
 A_p(y)=\int_0^\infty\phi(p+t)\phi(p-t)\cos(2yt)\,dt\quad(p\ge0).
\]
Changing variables $u=p+t,v=p-t$, and using evenness in $p,t$, gives
\[
 \Re\{F'(x+iy)\overline{F(x+iy)}\}
     =8\int_0^\infty p\sinh(2xp)A_p(y)\,dp.                 \tag{1}
\]
All interchanges are absolute before the bounded cosine is introduced.
The sufficient sublemma $A_p(y)\ge0$ for every $p>0,y$ would remove the
signed cancellation in (1). It is \emph{false for the actual theta kernel}.
It is not necessary for (1) to be positive.

\paragraph{Complete one-point certificate.}
The accompanying Arb script proves
\[
 -5.4320\cdot10^{-10}
 < A_0(20)=\int_0^\infty\phi(t)^2\cos(40t)\,dt
 < -5.4319\cdot10^{-10}.                                \tag{2}
\]
It is replayed at 256 and 384 bits with four theta terms and physical
cutoff $L=2$, then at 384 bits with eight terms and $L=4$.
This is one fixed transform value, not a scan, an off-line zeta zero,
a failure of RH, or a sign of the weighted integral (1).

Here are the complete error bounds. Every summand is positive for $t\ge0$.
Dropping its negative term bounds it above by
$2\pi^2n^4 e^{9t/2-\pi n^2e^{2t}}$. This expression decreases with $t$,
since its logarithmic derivative is $9/2-2\pi n^2e^{2t}<0$.
Writing $r=16e^{-3\pi}<1$, the decreasing successive ratios of
$n^4e^{-\pi n^2}$ give
\[
 0<\phi(t)\le C=\frac{2\pi^2e^{-\pi}}{1-r}<1,
 \quad
 0\le\phi(t)-\phi_M(t)\le B_M=
 \frac{2\pi^2(M+1)^4e^{-\pi(M+1)^2}}
 {1-[(M+2)/(M+1)]^4e^{-\pi(2M+3)}}.
\]
Consequently the replacement of $\phi^2$ by $\phi_M^2$ on $[0,L]$
costs at most $2LCB_M$. For $t\ge L$ we also have
\[
 \phi(t)\le C_0e^{9t/2-\pi e^{2t}},\qquad C_0=2\pi^2/(1-r).
\]
Using $e^{2t}\ge e^{2L}[1+2(t-L)]$ proves
\[
 \int_L^\infty\phi(t)^2dt
 \le\frac{C_0^2e^{9L-2\pi e^{2L}}}{4\pi e^{2L}-9},
\]
for the positive denominator at both selected cutoffs. The series error
at $(M,L)=(4,2)$ is below $3.277\cdot10^{-30}$ and the physical tail
below $3.923\cdot10^{-142}$. Arb integrates the entire finite prefix;
no midpoint solve or dropped exterior enters (2).

For completeness, $A_p(20)$ is continuous at $p=0$. On a bounded $t$
interval this follows from continuity and boundedness of the complete
kernel. For $0\le p\le1$ and $t\ge2$, both $|p\pm t|\ge t-1$.
The upper envelope $C_0e^{9u/2-\pi e^{2u}}$ is decreasing for $u\ge0$,
so its square at $u=t-1$ is an integrable common majorant for the
product. Dominated convergence applies. Thus (2) also implies
$A_p(20)<0$ for all sufficiently small positive $p$, without asserting
a numerical radius. The factor $p\sinh(2xp)$ in (1) does not rescue the
claim that \emph{each} slice is nonnegative.

\paragraph{Generic kernel convexity does not repair the argument.}
Let
\[
 k(t)=\frac1{2\sqrt\pi}\{3e^{-t^2/4}+e^{-(t-1)^2/4}+e^{-(t+1)^2/4}\}.
\]
This kernel is positive, smooth, even and strongly log-concave. Indeed
$\partial_t^2\log k=-1/2+\operatorname{Var}_t(c)/4\le-1/4$,
where the tilted centers are $c\in\{-1,0,1\}$ with positive weights
summing to one. Its bilateral Laplace transform is
$F_k(z)=e^{z^2}(3+2\cosh z)$. At $z=\log2+i\pi$,
\[
 \Re(F_k'\overline{F_k})
  =e^{2[(\log2)^2-\pi^2]}(\tfrac12\log2-\tfrac34)<0.
\]
Here $\log2<1$ is elementary. Its zeros include
$z=\pm\operatorname{arcosh}(3/2)+(2m+1)i\pi$, off the imaginary axis.
This is a control on positivity/evenness/ordinary log-concavity alone.
The model is order two and has no zeta arithmetic or theta modular law;
it does not refute a claim using those additional hypotheses. Nine
Maxima identities check the elementary algebra, not the classical
analytic criterion or the entire historical repository.

\paragraph{What would be a genuinely additional input?}
An explicit factorization of the \emph{integrated} kernel in (1) as an
autocorrelation, derived from the actual theta modular identity, would
supply information absent from the circle construction. Defining a
factor by taking the square root of the unknown Fourier transform assumes
the desired sign and is forbidden. Neither such a factor nor a signed
lower bound has been found in this reassessment. The per-slice factorization
is now excluded by (2). The broader integrated-kernel program is classical
and remains open; it must not be labeled a novel RH approach merely
because it is drawn in circle coordinates.
